% Page 608 — Thorne & Campolattaro (1967), ApJ 149, 591
% Appendix C: Equations of Motion for Even Parity

\appendix
\section*{APPENDIX C \\ EQUATIONS OF MOTION FOR EVEN PARITY}

The fluid 4-velocity corresponding to the even-parity equation of (7) in, to
first order in the displacement,
\begin{align}
u^{0} &= e^{-\nu/2}\bigl[1 - \tfrac{1}{2}H_{0}\,P_{l}(\cos\theta)\bigr] , &
u^{r} &= e^{-\lambda/2} r^{-1} e^{-(\lambda+\nu)/2} \dot{W}_{l}\,P_{l}(\cos\theta) ,
\nonumber\\
u^{\theta} &= -r^{-1} e^{-\nu/2} \dot{V}_{l}\,\partial_{\theta}P_{l}(\cos\theta) , &
u^{\phi} &= 0 \;.
\tag{C1}
\end{align}

In the pulsating configuration the Lagrangian change in the number density of baryons is
\begin{equation}
\Delta n = -n\,\Theta^{k}{}_{;k} - \tfrac{1}{2}n\,\dot{h}^{l}{}_{l}\,({}^{(3)}g)^{1/2} \;.
\tag{C2}
\end{equation}
Here $\Theta^{k}{}_{;k}$ is the divergence of the fluid displacement with respect to the 3-geometry at constant
time, $t$; and ${}^{(3)}g$ is the determinant of the metric of that 3-geometry.  In the Regge-Wheeler
choice of gauge, equation (C2) reduces to$^{3}$
\begin{equation}
\Delta n/n = \bigl[-r^{-2}e^{(\lambda-\nu)/2}\dot{W} - l(l+1)r^{-2}\dot{V}
             + \tfrac{1}{2}\dot{H}_{1} + \dot{K}\bigr]P_{l}(\cos\theta) \;.
\tag{C3}
\end{equation}

The corresponding Eulerian changes in pressure and in density of mass-energy are
\begin{align}
\delta p &= (p+\rho)\,(\Delta n/n) - p'\,r^{-1}e^{-(\lambda+\nu)/2}W\,P_{l}(\cos\theta) ,
\nonumber\\
\delta\rho &= \gamma p\,(\Delta n/n) - \rho'\,r^{-1}e^{-(\lambda+\nu)/2}W\,P_{l}(\cos\theta) \;;
\tag{C4}
\end{align}
and the Eulerian changes in the stress-energy tensor are
\begin{align}
\delta T^{0}{}_{0} &= -\delta\rho , &
\delta T^{1}{}_{1} &= \delta T^{2}{}_{2} = \delta T^{3}{}_{3} = \delta p , &
\delta T^{0}{}_{1} &= (p+\rho)\,u^{1} ,
\nonumber\\
\delta T^{1}{}_{0} &= -(p+\rho)\,u^{1} , &
\delta T^{2}{}_{0} &= (p+\rho)\,u^{0}\mu^{2} , &
\delta T^{3}{}_{0} &= (p+\rho)\,u^{0}\mu^{3} \;.
\tag{C5}
\end{align}

All other components of $\delta T^{\mu}{}_{\nu}$ vanish.

The way we had the perturbations in the Einstein field equations,
$\delta G^{\mu}{}_{\nu} = 8\pi\,\delta T^{\mu}{}_{\nu}$,
and in the equations of motion of the fluid,
$\delta(T^{\mu}{}_{\nu;\mu}) = 0$, associated with the perturbed
stress-energy tensor (C5) and the perturbed metric (7b); and we have checked our result on the
IBM 7094 using the computer programs of Thorne and Zimmerman (1967).$^{4}$
The origins of equations (8) and (9) are as follows.
Equation (8a) is $\delta G^{t}{}_{t} - \delta G^{r}{}_{r} = 8\pi\,\delta T^{t}{}_{t} - 8\pi\,\delta T^{r}{}_{r}$,
and it has been used to eliminate $H_{2}$ from all equations.
Equation (8b) is $\delta G^{r}{}_{t} = 8\pi\,\delta T^{r}{}_{t}$,
and it has been used to eliminate $H_{1}$ from all equations.
Equation (8c) is $\delta G^{r}{}_{r} = 8\pi\,\delta T^{r}{}_{r}$,
and it has been used to eliminate $\dot{K}$ from all other equations.
Equation (8d) is $\delta G^{\theta}{}_{\theta} = 8\pi\,\delta T^{\theta}{}_{\theta}$,
and it has been used to eliminate $\dot{H}_{1}$ from all other equations.
Equation (9a) is $\delta G^{r}{}_{\theta} = 0$.
Equation (9b) is $\delta(T^{r}{}_{\theta;\mu}) = 0$.
Combined with equation (8a), equation (9b) is $\delta G^{\theta}{}_{r} = 0$.
Equation (9c) is $\delta(T^{r}{}_{\mu;\nu}) = 0$.
Equation (9d) is $\delta(T^{r}{}_{\mu;\nu}) = 0$.
Several other useful dynamical equations for the perturbed configuration, which can be
derived from equations (2), (3), (8), and (9) with a large amount of effort, are these:
$\delta G^{r}{}_{t} = 8\pi\,\delta T^{r}{}_{t}$ is
\begin{equation}
H_{1}'' + \bigl[\tfrac{1}{2}r^{-1}(\lambda' - \nu') - 2r^{-1}\bigr]H_{1}'
  + \bigl[2r\nu' e^{-\nu} p - \nu'' - \tfrac{1}{4}(\nu')^{2}\bigr]H_{1}
  = e^{\lambda}(K_{1} + H_{00}) - 16\pi(p+\rho)V \;.
\tag{C6}
\end{equation}

\bigskip
\noindent\rule{0.4\textwidth}{0.4pt}

\footnotetext[3]{The peculiar form (7a) in which we chose to define $\delta n$ was motivated by a desire to take
on as simple a form as possible and to thereby simplify the eigensys.\ (14).}

\footnotetext[4]{$G^{r}{}_{r}$, $K_{r}$, $T^{r}{}_{r}$, and $T^{\phi}{}_{\phi}$ for even-parity pulsations are presented in Thorne and Zimmerman (1967)
precisely as they were output by the computer.}
